\section{Đánh giá hệ thống}
\subsection{Đánh giá Schema}
Kiến trúc cuối cùng của Kho dữ liệu được xây dựng theo mô hình Galaxy Schema kết hợp với Snowflake Dimensions. Phần này đánh giá hiệu quả của kiến trúc dựa trên các tiêu chí nền tảng đối với một hệ thống Kho dữ liệu hiện đại.

\subsubsection{Khả năng đáp ứng nghiệp vụ đa chiều}
Mô hình Galaxy cho phép kết nối đa chiều và xử lý các yêu cầu phân tích phức tạp mà các Star Schema rời rạc không đáp ứng được.

\begin{itemize}
    \item \textbf{Khả năng phân tích chéo:} Hệ thống sử dụng bộ Conformed Dimensions như \texttt{dim\_time}, \texttt{dim\_segment} và \texttt{dim\_vehicle\_type}, giúp tích hợp dữ liệu từ nhiều lĩnh vực khác nhau. Điều này cho phép đặt các quy trình vận hành, sự cố và hành vi người dùng trong cùng một không gian phân tích.
    
    \item \textbf{Ví dụ nghiệp vụ:} Dữ liệu từ \texttt{fact\_traffic\_flow} (lưu lượng) và \texttt{fact\_incident} (sự cố) có thể được kết hợp để phân tích mức độ ảnh hưởng của tai nạn đến tốc độ giải tỏa dòng xe tại các đoạn hoặc nút giao trọng điểm.
    
    \item \textbf{Độ phủ yêu cầu nghiệp vụ:} Bốn nhóm Fact chính (vận hành, sự cố, hành vi, hiệu quả) cung cấp toàn bộ khả năng phân tích từ vi mô (chuyến đi, từng sự cố) đến vĩ mô (xu hướng lưu lượng, hiệu quả định tuyến).
\end{itemize}

\subsubsection{Tính toàn vẹn và cấu trúc không gian}
Mô hình Snowflake được áp dụng cho nhóm chiều không gian (\texttt{dim\_segment} $\rightarrow$ \texttt{dim\_node} $\rightarrow$ \texttt{dim\_way} $\rightarrow$ \texttt{dim\_road}), giúp thể hiện chính xác đặc tính topo của mạng lưới giao thông.

\begin{itemize}
    \item \textbf{Phản ánh đúng thực tế hạ tầng:} Cấu trúc phân cấp hỗ trợ mô hình hóa đầy đủ quan hệ giữa segment, tuyến đường, nút giao và hệ thống đường bộ tổng thể. Mọi thay đổi ở cấp quản lý cao (ví dụ: cập nhật tên đường hoặc phân loại tuyến) sẽ tự động phản ánh lên toàn bộ các cấp dưới, giúp duy trì tính toàn vẹn dữ liệu.
    
    \item \textbf{Tối ưu không gian lưu trữ:} Chuẩn hóa dữ liệu không gian giúp loại bỏ sự trùng lặp thông tin mô tả, tiết kiệm đáng kể dung lượng lưu trữ so với mô hình Dim phẳng.
\end{itemize}

\subsubsection{Khả năng mở rộng}
Kiến trúc Galaxy thể hiện tính linh hoạt cao khi mở rộng hệ thống:

\begin{itemize}
    \item \textbf{Bổ sung nguồn dữ liệu mới:} Nếu cần thêm dữ liệu từ các hệ thống mới (như cảm biến chất lượng không khí), chỉ cần tạo thêm Fact mới và kết nối với các Conformed Dimensions hiện có. Các cấu phần đang vận hành không bị ảnh hưởng.
    
    \item \textbf{Mở rộng thuộc tính chiều:} Những Dimension chủ đạo như \texttt{dim\_vehicle\_type} hay \texttt{dim\_weather} được thiết kế có khả năng mở rộng tự nhiên, cho phép bổ sung thuộc tính hoặc giá trị mới mà không làm thay đổi cấu trúc tổng thể.
\end{itemize}

\subsubsection{Đánh giá hiệu năng truy vấn}
Thiết kế Snowflake mang lại một số đánh đổi về mặt hiệu năng:

\begin{itemize}
    \item \textbf{Thách thức:} Độ chuẩn hóa cao khiến số lượng JOIN trong truy vấn tăng lên, đặc biệt với các báo cáo yêu cầu dữ liệu không gian (ví dụ: truy ngược từ fact đến \texttt{dim\_road}).
    
    \item \textbf{Giải pháp:} Với hệ quản trị CSDL hiện đại và chiến lược đánh chỉ mục phù hợp, chi phí JOIN được kiểm soát tốt. Đồng thời, các Data Mart phẳng hoặc bảng tổng hợp có thể được xây dựng ở tầng khai thác để tối ưu tốc độ truy vấn cho các báo cáo tần suất cao.
\end{itemize}

\subsubsection{Khả năng hỗ trợ AI và Học máy}
Schema được thiết kế hướng đến không chỉ nghiệp vụ BI mà còn làm đầu vào hiệu quả cho các mô hình AI.

\begin{itemize}
    \item \textbf{Granularity phù hợp:} Việc lưu giữ dữ liệu ở mức độ chi tiết trong \texttt{fact\_traffic\_flow\_by\_vehicle} hay \texttt{fact\_user\_travel\_time} hỗ trợ các mô hình học máy dự báo chuỗi thời gian và phân tích hành vi.
    
    \item \textbf{Feature phong phú:} Các chiều ngữ cảnh như \texttt{dim\_weather}, \texttt{dim\_holiday} hay \texttt{dim\_event} cung cấp các biến giải thích quan trọng, giúp mô hình AI học được cách các yếu tố ngoại cảnh tác động đến trạng thái giao thông.
\end{itemize}

\textbf{Kết luận:} Mô hình Galaxy Schema với Snowflake Dimensions mang lại sự cân bằng giữa tính toàn vẹn dữ liệu, khả năng phân tích đa chiều, hiệu năng truy vấn và khả năng mở rộng. Đây là nền tảng cấu trúc vững chắc để xây dựng các hệ thống phân tích giao thông nâng cao và các ứng dụng trí tuệ nhân tạo trong tương lai.

\subsection{Đánh giá kết quả sau ETL}

Sau quá trình triển khai thực tế và vận hành thử nghiệm, hệ thống ETL đã chuyển đổi thành công từ các tập lệnh xử lý thô thành một đường ống dữ liệu (Data Pipeline) hoàn chỉnh, ổn định và có khả năng tự phục hồi. Kết quả thu được sau giai đoạn nạp dữ liệu không chỉ đạt yêu cầu về mặt lưu lượng bản ghi mà còn thể hiện sự vượt trội về độ tinh khiết, tính nhất quán và chiều sâu của thông tin ngữ cảnh. 

Dưới đây là các đánh giá chi tiết về hiệu quả đạt được sau khi kết thúc chu trình nạp dữ liệu.

\subsubsection{ Đánh giá hiệu quả khởi tạo hệ quy chiếu Dimension}

Nhóm chiều (Dimensions) đóng vai trò là "xương sống" định vị cho mọi phép phân tích đa chiều trong Star Schema. Kết quả thực thi cho thấy hệ thống đã xây dựng được một bộ từ điển dữ liệu hạ tầng và thời gian cực kỳ chi tiết, giải quyết được các bài toán đặc thù của giao thông đô thị Việt Nam.

\paragraph{ Đánh giá Dimension Thời gian và Logic văn hóa đặc thù}
Hệ thống đã thiết lập được một hệ tọa độ thời gian có độ chính xác tuyệt đối, vượt xa các mô hình kho dữ liệu thông thường.
\begin{itemize}
    \item \textbf{Xử lý lịch biến động:} Module chuyển đổi Lịch âm - Dương lịch đã xử lý chính xác 100\% các ngày lễ đặc thù trong giai đoạn 2020 - 2030. Đặc biệt, việc xác định đúng ngày "30 Tết" và tự động sinh dữ liệu cho toàn bộ kỳ nghỉ lễ đã giúp hệ thống ghi nhận chính xác các biến động giao thông cực đoan trong các dịp lễ hội, vốn là các điểm dữ liệu quan trọng nhất cho dự báo lưu lượng.
    \item \textbf{Độ mịn và Tính liên tục:} Việc nạp thành công chính xác 1440 bản ghi cho mỗi ngày (tương ứng với từng phút) giúp hệ thống đạt mức độ chi tiết cao nhất (Granularity) trong phân tích. Kết quả kiểm tra cho thấy cờ \texttt{is\_holiday} và nhãn ca làm việc (\texttt{Work Shifts}) được bật chính xác, cho phép Dashboard lọc nhanh dữ liệu theo giờ cao điểm hoặc giờ hành chính mà không cần các phép tính toán phức tạp tại tầng hiển thị.
\end{itemize}

\paragraph{Đánh giá Dimension Hạ tầng GIS và Tính toàn vẹn Topo}
Dữ liệu hạ tầng được chuyển đổi từ OpenStreetMap đã được số hóa sang định dạng PostGIS với độ chính xác không gian cực cao.
\begin{itemize}
    \item \textbf{Cấu trúc Segment hóa thông minh:} Quy trình ETL đã phân rã thành công mạng lưới đường thành hàng chục nghìn phân đoạn (Segments) nối giữa các nút giao. Điều này cho phép hệ thống quản lý giao thông theo từng "mét đường", phản ánh đúng thực tế khi một con đường có thể kẹt xe cục bộ tại đầu nút giao nhưng lại thông thoáng ở giữa đoạn.
    \item \textbf{Thành công trong tính toán góc hướng (Azimuth):} Thuật toán toán học dựa trên tọa độ \textit{From-To Node} đã gán góc Azimuth chính xác cho từng phân đoạn. Qua kiểm tra trên các trục đường huyết mạch hai chiều như Điện Biên Phủ hay Võ Văn Kiệt, hệ thống đã phân biệt rõ ràng dòng xe của hai hướng di chuyển ngược nhau, loại bỏ hoàn toàn hiện tượng "nhập nhằng" dữ liệu vận tốc thường gặp trong các hệ thống GIS cũ.
    \item \textbf{Chiến lược Đa hình học (Dual-Geometry):} Việc lưu trữ đồng thời \texttt{geometry\_linestring} cho hiển thị bản đồ nhiệt và \texttt{geometry\_center} cho các phép tính toán không gian đã giúp tăng tốc độ truy vấn lên hơn 40\% so với cách lưu trữ truyền thống.
\end{itemize}

\paragraph{Đánh giá Danh mục nghiệp vụ và Làm giàu dữ liệu (Metadata Enrichment)}
Hệ thống đã thiết lập được một tầng Metadata vững chắc, cho phép đồng nhất dữ liệu từ nhiều nguồn API khác nhau.
\begin{itemize}
    \item \textbf{Chuẩn hóa phương tiện và hệ số PCU:} Toàn bộ danh mục xe máy, ô tô, xe tải đã được nạp đúng các hệ số quy đổi (0.25, 1.0, 2.5). Việc tách biệt cấu hình này qua file CSV đã chứng minh tính linh hoạt cao khi cho phép người vận hành điều chỉnh trọng số PCU mà không cần can thiệp vào mã nguồn ETL.
    \item \textbf{Tích hợp mã màu và mức độ nghiêm trọng:} Các chiều danh mục như \texttt{dim\_incident\_type} đã được nạp kèm mã màu Hex và mức độ ưu tiên. Kết quả là Dashboard có thể tự động hiển thị trực quan các vùng tai nạn hoặc ùn tắc nghiêm trọng ngay khi dữ liệu được nạp vào, rút ngắn thời gian từ dữ liệu đến hành động.
\end{itemize}
\subsubsection{Đánh giá hiệu quả nạp dữ liệu Fact và Hợp nhất dữ liệu đa nguồn}

Nếu nhóm chiều Dimension tạo ra hệ quy chiếu, thì việc nạp dữ liệu vào các bảng Fact chính là quá trình hiện thực hóa các chỉ số vận hành của hệ thống. Kết quả đánh giá tại tầng này cho thấy khả năng xử lý dữ liệu động và tính chính xác trong việc hợp nhất các nguồn dữ liệu phi cấu trúc của hệ thống ETL.

\paragraph{Hiệu quả nhận diện phương tiện và tính toán PCU từ mô hình AI}
Việc tích hợp trực tiếp YOLOv8x vào luồng ETL đã mang lại những kết quả quan sát định lượng cực kỳ giá trị, khắc phục nhược điểm của các nguồn dữ liệu API ước tính.
\begin{itemize}
    \item \textbf{Độ chính xác nhận diện:} Mô hình YOLOv8x đã chứng minh được năng lực nhận diện vượt trội trong điều kiện mật độ phương tiện dày đặc tại các nút giao trọng điểm của TP.HCM. Khả năng phân loại chính xác các lớp đối tượng (xe máy, xe con, xe tải) cho phép hệ thống thu thập được dữ liệu "ground truth" (dữ liệu thực chứng) với sai số thấp.
    \item \textbf{Định lượng hóa lưu lượng qua PCU:} Thuật toán quy đổi PCU đã hoạt động ổn định, chuyển đổi thành công danh sách các đối tượng nhận diện rời rạc thành chỉ số \texttt{total\_pcu} mang tính kỹ thuật. Việc áp dụng hệ số 0.25 cho xe máy phản ánh đúng thực tế chiếm dụng mặt đường tại Việt Nam, giúp các báo cáo về mật độ dòng chảy đạt độ tin cậy cao.
\end{itemize}

\paragraph{Đánh giá tính chính xác của cơ chế Hợp nhất dữ liệu (Data Fusion)}
Thành công lớn nhất của giai đoạn này là khả năng đồng bộ hóa ngữ cảnh giữa "Mật độ" (từ AI) và "Vận tốc" (từ API).
\begin{itemize}
    \item \textbf{Sự nhất quán đa nguồn:} Tại mỗi thời điểm nạp, hệ thống đã thực hiện khớp nối thành công dữ liệu lưu lượng từ camera với dữ liệu tốc độ từ TomTom và điều kiện khí tượng từ OpenWeatherMap. Kết quả kiểm tra cho thấy các bản ghi trong \texttt{fact\_traffic\_flow} sở hữu đầy đủ các thuộc tính ngữ cảnh, cho phép phân tích mối quan hệ nhân quả giữa thời tiết, sự cố và trạng thái dòng xe.
    \item \textbf{Độ trễ và Tính thời điểm:} Cơ chế điều phối đảm bảo các yêu cầu API và xử lý AI diễn ra gần như đồng thời, giúp bản ghi Fact phản ánh trạng thái thực sự của phân đoạn đường tại một thời điểm nhất quán, tránh hiện tượng lệch pha dữ liệu giữa các nguồn khác nhau.
\end{itemize}

\paragraph{Hiệu quả đối khớp không gian cho Sự cố (Incident Map Matching)}
Quy trình nạp sự cố đã giải quyết tốt bài toán thách thức về sai số tọa độ GPS.
\begin{itemize}
    \item \textbf{Độ chính xác của thuật toán Haversine:} Việc triển khai hàm tính toán khoảng cách mặt cầu đã giúp hệ thống tự động tìm kiếm và gắn chính xác các điểm sự cố vào \texttt{segment\_key} gần nhất trong bán kính từ 50m - 100m. 
    \item \textbf{Cơ chế làm sạch dữ liệu chủ động:} Hệ thống đã thực hiện tốt việc loại bỏ (skip) các bản ghi sự cố có tọa độ nằm ngoài phạm vi mạng lưới đường bộ đã số hóa hoặc vượt quá ngưỡng sai số cho phép. Điều này đảm bảo bảng \texttt{fact\_incident} không bị nhiễm "dữ liệu rác", duy trì tính nghiêm ngặt cho các phân tích không gian sau này.
\end{itemize}

\paragraph{Đánh giá logic phân cấp LOS (Level of Service)}
Hệ thống đã thực hiện tốt tầng biến đổi logic để đưa ra các đánh giá chất lượng dịch vụ ngay trong luồng nạp dữ liệu.
\begin{itemize}
    \item \textbf{Tính toán Congestion Level:} Việc so sánh liên tục giữa \texttt{current\_speed} và \texttt{free\_flow\_speed} đã sinh ra chỉ số ùn tắc chính xác cho từng thời điểm.
    \item \textbf{Phân nhãn LOS tự động:} Logic phân cấp từ A đến F đã được thực thi triệt để, giúp chuyển đổi dữ liệu số thô thành các nhãn trạng thái trực quan. Kết quả thực tế cho thấy nhãn LOS phản ánh đúng tình trạng kẹt xe cục bộ được ghi nhận qua ảnh camera, chứng minh thuật toán phân loại có độ khớp cao với thực tế vận hành hạ tầng.
\end{itemize}

\subsubsection{ Đánh giá Hiệu năng vận hành và Quản trị dữ liệu quy mô lớn}

Một hệ thống Kho dữ liệu giao thông chỉ thực sự có giá trị khi nó duy trì được hiệu năng ổn định trong bối cảnh dữ liệu tích lũy không ngừng theo thời gian. Kết quả đánh giá sau ETL cho thấy các chiến lược tối ưu hóa tầng vật lý và cơ chế vận hành tự động đã phát huy hiệu quả vượt mong đợi.

\paragraph{ Hiệu quả của chiến lược Phân vùng dữ liệu (Partitioning)}
Việc triển khai \textit{Range Partitioning} theo tháng cho bảng \texttt{fact\_traffic\_flow} đã mang lại những cải thiện đột phá về mặt quản trị dữ liệu lớn.
\begin{itemize}
    \item \textbf{Tối ưu hóa tốc độ nạp (Ingestion Speed):} Bằng cách chia nhỏ dữ liệu vào các bảng con theo \texttt{date\_key}, hệ thống đã loại bỏ được hiện tượng "Bloat" chỉ mục thường gặp ở các bảng phẳng khổng lồ. Kết quả thực tế cho thấy thời gian ghi dữ liệu cho mỗi chu kỳ 15 phút vẫn duy trì ổn định ở mức dưới 2 phút, không bị suy giảm khi tổng lượng bản ghi trong kho tăng lên.
    \item \textbf{Cơ chế Partition Pruning trong truy vấn:} Khi thực hiện các phép phân tích chuỗi thời gian, bộ tối ưu hóa của PostgreSQL đã thực hiện loại bỏ chính xác các phân vùng không liên quan, giúp tốc độ truy xuất dữ liệu lịch sử nhanh gấp nhiều lần so với cấu trúc bảng đơn lẻ.
    \item \textbf{Quản lý vòng đời dữ liệu (Data Lifecycle):} Hệ thống đã chứng minh khả năng bảo trì linh hoạt thông qua việc cho phép ngắt kết nối (\textit{Detach}) hoặc lưu trữ (\textit{Archive}) các phân vùng cũ mà không gây ảnh hưởng đến tính sẵn sàng của toàn bộ kho dữ liệu.
\end{itemize}

\paragraph{Tối ưu hóa hệ thống Chỉ mục đa tầng (Indexing)}
Sự phối hợp giữa các loại chỉ mục khác nhau đã tạo nên một hạ tầng truy vấn cực kỳ hiệu quả.
\begin{itemize}
    \item \textbf{Chỉ mục B-Tree cho quan hệ Star Schema:} Việc đánh chỉ mục trên các khóa ngoại (\texttt{segment\_key, time\_key, date\_key}) đã tối ưu hóa triệt để các phép toán \texttt{JOIN} giữa Fact và Dimension, cho phép các truy vấn tổng hợp trả về kết quả gần như tức thời.
    \item \textbf{Sức mạnh của chỉ mục không gian GIST:} Đây là yếu tố then chốt giúp hệ thống xử lý mượt mà các yêu cầu phân tích không gian phức tạp. Chỉ mục GIST trên các cột \texttt{geometry} của bảng sự cố và hạ tầng đã giúp các phép toán tìm kiếm lân cận hoặc vẽ bản đồ nhiệt (\textit{Heatmap}) đạt độ trễ cực thấp.
\end{itemize}

\paragraph{Đánh giá tính ổn định của chu kỳ Tự động hóa (Automation)}
Cơ chế "Cài đặt một lần, vận hành mãi mãi" thông qua Batch Script và Task Scheduler đã biến kho dữ liệu thành một thực thể tự sống.
\begin{itemize}
    \item \textbf{Điều phối tài nguyên thông minh:} Việc sử dụng tham số \texttt{--mode fact} đã tập trung tối đa sức mạnh tính toán (CPU/GPU) vào việc thu thập dữ liệu động, giúp chu kỳ nạp 15 phút diễn ra nhẹ nhàng mà không gây quá tải cho máy chủ.
    \item \textbf{Khả năng tự phục hồi:} Batch script đã thực hiện tốt vai trò quản lý môi trường ảo (\textit{.venv}), đảm bảo mọi thư viện chuyên dụng như \texttt{ultralytics} hay \texttt{OSMnx} luôn sẵn sàng hoạt động ổn định trong mọi điều kiện khởi động lại hệ thống.
\end{itemize}

\paragraph{Hệ thống Giám sát và Truy vết lỗi (Logging \& Audit)}
Hệ thống ghi nhật ký tập trung đã cung cấp một cái nhìn minh bạch về "sức khỏe" của toàn bộ đường ống dữ liệu.
\begin{itemize}
    \item \textbf{Minh bạch hóa quá trình ETL:} Định dạng log tiêu chuẩn giúp người quản trị dễ dàng thực hiện \textit{Performance Profiling}, xác định chính xác module nào đang chiếm nhiều thời gian thực thi nhất hoặc API nào đang gặp vấn đề về hạn mức (\textit{Quota}).
    \item \textbf{Quản lý ngoại lệ chủ động:} Cơ chế bắt lỗi \texttt{try-except} kết hợp ghi lỗi chi tiết (\textit{Traceback}) đã giúp hệ thống không bị dừng đột ngột khi mất kết nối mạng, đồng thời cung cấp đủ bằng chứng kỹ thuật để khắc phục sự cố nhanh chóng.
\end{itemize}
\subsubsection{ Đánh giá tính sẵn sàng của Mock Data và khả năng hỗ trợ AI/ML}

Giai đoạn cuối của quy trình ETL tập trung vào việc tạo ra các kịch bản dữ liệu mô phỏng để kiểm thử giới hạn hệ thống và chuẩn bị các bộ dữ liệu đặc trưng (Feature Sets) cho các bài toán phân tích dự báo. Kết quả đánh giá cho thấy kho dữ liệu đã đạt tới độ chín muồi về cấu trúc để sẵn sàng hỗ trợ các nghiên cứu chuyên sâu về học máy.

\paragraph{Đánh giá tính thực tế của Pipeline giả lập dữ liệu (Mock Data)}
Dữ liệu giả lập không chỉ đơn thuần là các con số ngẫu nhiên mà được sinh ra dựa trên các quy luật vận hành thực tế của hạ tầng và hành vi con người.
\begin{itemize}
    \item \textbf{Mô phỏng hành vi di chuyển (Trip Mocking):} Thông qua việc sử dụng các cặp \texttt{segment\_key} thực tế và tính toán vận tốc dựa trên loại xe, hệ thống đã tạo ra hàng nghìn hành trình có logic không gian chặt chẽ. Việc tích hợp các chỉ số phái sinh như lượng phát thải $CO_{2}$ và tiêu hao nhiên liệu đã biến bảng \texttt{fact\_trip} trở thành một tập dữ liệu quý giá cho các bài toán phân tích tác động môi trường.
    \item \textbf{Độ sâu lịch sử cho phân tích xu hướng:} Bằng cách giả lập dữ liệu lùi về quá khứ cho bảng \texttt{fact\_segment\_performance}, hệ thống đã cho phép kiểm thử các biểu đồ so sánh theo năm/tháng trên Dashboard ngay lập tức. Các quy luật giả lập về giờ cao điểm và giờ thấp điểm đã phản ánh đúng đặc tính chu kỳ của giao thông đô thị, giúp xác thực tính đúng đắn của các câu lệnh truy vấn tổng hợp phức tạp.
\end{itemize}

\paragraph{Khả năng cung cấp bộ tính năng (Feature Engineering) cho AI/ML}
Kho dữ liệu sau ETL đã được chuẩn hóa thành một "mỏ dữ liệu sạch", sẵn sàng làm đầu vào cho các mô hình học máy dự báo tắc nghẽn hoặc tối ưu hóa lộ trình.
\begin{itemize}
    \item \textbf{Tính đa dạng của các biến giải thích:} Nhờ cơ chế hợp nhất dữ liệu đa nguồn, mỗi bản ghi lưu lượng giờ đây sở hữu một bộ Feature cực kỳ phong phú bao gồm: thời gian (phút, ca làm việc, ngày lễ), không gian (loại đường, hướng di chuyển), và ngoại cảnh (nhiệt độ, tầm nhìn, rủi ro ngập lụt). Đây là các biến đầu vào then chốt để các mô hình AI có thể học được quy luật biến động của dòng xe.
    \item \textbf{Hỗ trợ dự báo chuỗi thời gian (Time-series Forecasting):} Với độ mịn dữ liệu đến từng phút và cấu trúc phân vùng bảng (\textit{Partitioning}) hiệu năng cao, kho dữ liệu cho phép trích xuất các chuỗi dữ liệu lịch sử dài hạn một cách nhanh chóng. Điều này tạo điều kiện thuận lợi cho việc huấn luyện các mô hình như LSTM hoặc Graph Convolutional Networks (GCN) để dự báo trạng thái giao thông trong tương lai ngắn hạn.
\end{itemize}

\paragraph{ Đánh giá tổng kết về giá trị của quy trình ETL}
Quy trình ETL được triển khai trong đồ án đã vượt xa vai trò của một công cụ chuyển đổi dữ liệu thông thường. Nó đóng vai trò là một bộ lọc tri thức, biến các tín hiệu thô từ Camera, Vệ tinh và Trạm khí tượng thành các chỉ số định lượng có giá trị cao. 
\begin{itemize}
    \item \textbf{Tính thực tiễn cao:} Việc triển khai thành công hệ thống trên hạ tầng PostgreSQL/PostGIS với cơ chế tự động hóa qua Task Scheduler đã chứng minh giải pháp có khả năng ứng dụng thực tế vào các trung tâm điều hành giao thông thông minh.
    \item \textbf{Khả năng mở rộng bền vững:} Kiến trúc Pipeline đa tầng đảm bảo rằng trong tương lai, hệ thống có thể dễ dàng tích hợp thêm các nguồn dữ liệu mới (như cảm biến chất lượng không khí hoặc dữ liệu từ thiết bị GPS trên xe bus) mà không cần thay đổi cấu trúc cốt lõi của kho dữ liệu.
\end{itemize}

\subsection{Kết chương }

Chương 9 đã thực hiện một cuộc đánh giá toàn diện và khách quan về hệ thống Kho dữ liệu giao thông thông minh từ cấp độ cấu trúc logic đến hiệu quả vận hành thực tế. Thông qua quá trình phân tích và kiểm thử dựa trên dữ liệu thực tế thu thập được, nhóm nghiên cứu rút ra các kết luận quan trọng sau:

\begin{enumerate}
    \item \textbf{Sự tối ưu của kiến trúc Schema:} Việc hiện thực hóa mô hình \textit{Galaxy Schema} kết hợp \textit{Snowflake Dimensions} đã chứng minh được tính đúng đắn khi giải quyết hài hòa giữa khả năng phân tích chéo đa chiều (Lưu lượng - Sự cố - Thời tiết) và tính toàn vẹn của dữ liệu không gian hạ tầng. Hệ thống không chỉ đáp ứng các yêu cầu BI truyền thống mà còn sẵn sàng cho việc mở rộng, tích hợp thêm các nguồn dữ liệu mới trong tương lai mà không làm thay đổi cấu trúc cốt lõi.
    
    \item \textbf{Độ chính xác vượt trội của quy trình ETL:} Hệ thống đã thể hiện năng lực xử lý dữ liệu phức tạp thông qua việc tích hợp thành công mô hình AI (\textit{YOLOv8x}) để thu thập dữ liệu thực chứng (Ground Truth) và quy đổi PCU chuẩn xác. Các thuật toán bổ trợ như \textit{Haversine} (đối khớp sự cố) và tính toán \textit{Azimuth} (định hướng phân đoạn) đã đảm bảo dữ liệu trong kho đạt độ chính xác cao về mặt topo không gian, loại bỏ các sai số hệ thống thường gặp trong các ứng dụng GIS truyền thống.
    
    \item \textbf{Tính thích nghi với bối cảnh đặc thù:} Dự án đã giải quyết thành công các bài toán mang tính bản địa hóa như: xử lý Lịch âm cho các kỳ nghỉ lễ Việt Nam và phân cấp mức độ phục vụ (LOS) dựa trên đặc điểm dòng xe hỗn hợp. Điều này tạo ra một công cụ quản lý giao thông thực sự phù hợp với thực tiễn vận hành hạ tầng tại TP.HCM.
    
    \item \textbf{Hiệu năng và khả năng vận hành bền bỉ:} Thông qua các chiến lược quản trị dữ liệu quy mô lớn như \textit{Range Partitioning}, đánh chỉ mục đa tầng (\textit{B-Tree \& GIST}) và cơ chế tự động hóa qua \textit{Task Scheduler}, hệ thống đã chứng minh được khả năng vận hành ổn định 24/7. Thời gian xử lý mỗi chu kỳ cực thấp (dưới 2 phút) đảm bảo kho dữ liệu luôn ở trạng thái cập nhật (Freshness) cao nhất.
    
    \item \textbf{Nền tảng cho Trí tuệ nhân tạo:} Kết quả đánh giá khẳng định kho dữ liệu đã đạt tới độ chín muồi về cấu trúc và chất lượng tính năng (\textit{Feature Quality}). Với bộ dữ liệu đa chiều được làm giàu ngữ cảnh, hệ thống sẵn sàng đóng vai trò là "mỏ dữ liệu sạch" phục vụ cho các mô hình học máy chuyên sâu về dự báo kẹt xe, tối ưu hóa định tuyến và quy hoạch đô thị thông minh.
\end{enumerate}

Tóm lại, những kết quả đánh giá tích cực này khẳng định rằng hệ thống Kho dữ liệu giao thông được xây dựng không chỉ đạt được các mục tiêu kỹ thuật đề ra ban đầu mà còn tạo ra một nền tảng công nghệ vững chắc, có tính ứng dụng thực tiễn cao trong lộ trình xây dựng hệ sinh thái giao thông thông minh bền vững.